% Partie : méthodes
% Chapitre : induction mathématique
% Section : induction structurelle

\documentclass[../../../include/open-logic-section]{subfiles}

\begin{document}
\olfileid[fr]{mth}{ind}{sti}

\olsection{Induction structurelle}

Jusqu'ici, nous avons
utilisé la récurrence
pour démontrer des
résultats concernant
tous les nombres
naturels. Un principe
semblable s'applique
directement à tous
les !!{element}s
d'un ensemble défini
inductivement. On
l'appelle souvent
induction
\emph{structurelle},
car il s'appuie
sur la structure
des objets définis.

En général, une
définition inductive
donne (a)~une liste
d'!!{element}s
«~initiaux~» de
l'ensemble et
(b)~une liste
d'opérations qui
produisent de
nouveaux
!!{element}s
à partir de
ceux déjà construits.
Pour les termes
bien formés,
les objets initiaux
sont les lettres
et il n'y a
qu'une opération :
\begin{align*}
  o(s_1,s_2)
    &= [s_1 \circ s_2].
\end{align*}
On peut même voir
les nombres naturels
$\Nat$ comme un
ensemble défini
inductivement :
l'objet initial
est~$0$, et
l'opération est
la fonction
successeur $x+1$.

Pour démontrer que
tous les
!!{element}s
d'un ensemble
défini inductivement
possèdent une
propriété~$P$,
il faut :
\begin{enumerate}
\item montrer que
  les objets initiaux
  possèdent~$P$ ;
\item montrer que,
  pour chaque
  opération~$o$,
  si ses arguments
  possèdent~$P$,
  son résultat
  la possède aussi.
\end{enumerate}
Ainsi, pour démontrer
une propriété de
tous les termes
bien formés,
on la démontre
pour chaque lettre,
puis pour
$[s_1 \circ s_2]$
en supposant
qu'elle vaut
séparément pour
$s_1$ et~$s_2$.

\begin{prop}
  Dans tout terme
  bien formé~$t$,
  le nombre de
  crochets ouvrants
  $[$ est égal
  au nombre de
  crochets fermants
  $]$.
\end{prop}

\begin{proof}
Procédons par
induction structurelle.
Les termes bien
formés sont définis
à partir des lettres
au moyen de
l'opération~$o$.
\begin{enumerate}
\item Pour chaque
  lettre, la
  propriété est
  vraie : elle
  contient $0$
  crochet ouvrant
  et $0$ crochet
  fermant.
\item Supposons
  que les nombres
  de $[$ et de
  $]$ soient égaux
  dans~$s_1$,
  et qu'il en
  aille de même
  dans~$s_2$.
  Dans
  $o(s_1,s_2)
  = [s_1 \circ s_2]$,
  le nombre de
  $[$ est la
  somme des nombres
  de $[$ dans
  $s_1$ et $s_2$,
  plus~$1$.
  Le nombre de
  $]$ est la
  somme des nombres
  de $]$ dans
  $s_1$ et $s_2$,
  plus~$1$.
  Ils sont donc
  égaux dans
  $o(s_1,s_2)$.
\end{enumerate}
\end{proof}

\begin{prob}
  Démontrez par
  induction structurelle
  qu'aucun terme
  bien formé ne
  commence par~$]$.
\end{prob}

Voici une autre
démonstration par
induction structurelle.
Un \emph{préfixe propre
non vide} d'une
suite de symboles~$t$
est une suite~$s$
non vide qui coïncide
avec le début
de~$t$, symbole
par symbole, mais
qui est plus courte
que~$t$. Ainsi,
$[a \circ {}$
est un préfixe
propre non vide
de $[a \circ b]$ ;
en revanche,
$[b \circ {}$
ne l'est pas
(les deux suites
diffèrent au
deuxième symbole),
et $[a \circ b]$
non plus (il
a la même
longueur).

\begin{prop}\ollabel{prop:initial}
  Tout préfixe
  propre non vide
  d'un terme
  bien formé~$t$
contient plus
de $[$ que
de $]$.\footnote{
Correction de l'édition française :
la source anglaise définit un
segment initial propre sans
exclure la suite vide, puis
affirme qu'il contient plus
de crochets ouvrants que de
crochets fermants. La suite
vide est pourtant un préfixe
propre de tout terme non vide
et contient zéro crochet de
chaque sorte. Le mot
«~non vide~» est donc
nécessaire à l'énoncé
et à sa démonstration.}
\end{prop}

\begin{proof}
Raisonnons par
induction sur~$t$.
\begin{enumerate}
\item Si $t$
  est une lettre,
  il n'a aucun
  préfixe propre
  non vide.
\item Supposons
  $t=[s_1 \circ s_2]$
  pour des termes
  bien formés
  $s_1$ et~$s_2$.
  Si $r$ est
  un préfixe
  propre non vide
  de~$t$,
  plusieurs cas
  sont possibles :
  \begin{enumerate}
  \item $r$ est
    simplement $[$ :
    il contient
    un $[$ de
    plus que
    de $]$.
  \item $r$ est
    $[r_1$, où
    $r_1$ est
    un préfixe
    propre non
    vide de~$s_1$.
    Par hypothèse
    d'induction,
    $r_1$ contient
    plus de $[$
    que de $]$.
    Il en va
    donc de même
    pour $[r_1$.
  \item $r$ est
    $[s_1$ ou
    $[s_1 \circ {}$.
    D'après la
    proposition
    précédente,
    les nombres
    de $[$ et
    de $]$ dans
    $s_1$ sont
    égaux. Dans
    chacun de ces
    deux préfixes,
    $[s_1$ ou
    $[s_1 \circ {}$,
    le $[$ extérieur
    en ajoute un.
  \item $r$ est
    $[s_1 \circ r_2$,
    où $r_2$ est
    un préfixe
    propre non
    vide de~$s_2$.
    Par hypothèse
    d'induction,
    $r_2$ contient
    plus de $[$
    que de $]$.
    La proposition
    précédente donne
    autant de
    crochets de
    chaque sorte
    dans~$s_1$.
    Le $[$ extérieur
    rend donc
    l'inégalité
    stricte pour
    $[s_1 \circ r_2$.
  \item $r$ est
    $[s_1 \circ s_2$.
    D'après la
    proposition
    précédente,
    $s_1$ et
    $s_2$ ont
    chacun autant
    de $[$ que
    de $]$.
    Le préfixe
    $[s_1 \circ s_2$
    contient donc
    un $[$ de plus.
  \end{enumerate}
\end{enumerate}
\end{proof}

\end{document}
